\begin{abstract}
Spiking Neural Networks (SNNs) are frequently motivated by an appeal to biological
realism, on the assumption that mechanisms closer to their cortical counterparts will
perform better. We treat that assumption as a hypothesis to be measured rather than a
design principle to be followed. We introduce a spatially aware connectivity scheme in
which input-to-excitatory projections are elongated, orientation-tuned two-dimensional
Gaussian \textit{receptive fields} (RFs), and combine it with supervised reward-modulated
spike-timing-dependent plasticity (R-STDP), which assigns each excitatory neuron to a
class and reshapes its field through a baseline-corrected teaching signal. Varying the
connectivity, the learning rule, and the dataset independently, we find that the
structural prior contributes conditionally, and that the rule acting on it decides how
much. Under unsupervised plasticity the prior is worth between $0.6$ and $3.5$ percentage
points on all five datasets, yet no unsupervised rule improves on simply freezing the
weights, and the rules differ sharply in how much of the prior they leave standing, from
$93\%$ of the initial orientation coherence under triplet-STDP to $54\%$ under
trace-STDP. Under R-STDP the same contrast grows to $+4.9$ points on MNIST and $+5.3$ on
KMNIST, and reverses to $-2.6$ on notMNIST, whose class-relevant information does not
reside in local oriented structure. The model reaches $94.6\%$ on MNIST through its own
online readout, and its predictive entropy separates correct from incorrect predictions
well enough to reach $98.1\%$ accuracy at $90\%$ coverage---but only where the readout is
already accurate, the separation falling from $\mathrm{AUROC}$ $0.933$ on MNIST to
$0.595$ on SVHN. Biological realism is therefore better understood not as a scalar to be
maximized, but as an alignment between a structural assumption, the plasticity rule that
acts on it, and the statistical structure of the data.
\end{abstract}
